\documentclass[12pt]{article}
\usepackage{stmaryrd}
\usepackage{graphicx} % Pour l'insertion d'images
\usepackage{float}    % Pour contrôler précisément le placement
\usepackage[utf8]{inputenc}

\usepackage[french]{babel}
\usepackage[T1]{fontenc}
\usepackage{hyperref}
\usepackage{verbatim}

\usepackage{color, soul}

\usepackage{pgfplots}
\pgfplotsset{compat=1.15}
\usepackage{mathrsfs}

\usepackage{amsmath}
\usepackage{amsfonts}
\usepackage{amssymb}
\usepackage{tkz-tab}
\author{Destiné aux élèves de la 1er S2\\Lycée de Dindéfelo\\Présenté par M. BA}
\title{\textbf{Equations-Inéquations-Systèmes}}
\date{\today}
\usepackage{tikz}
\usetikzlibrary{arrows, shapes.geometric, fit}

% Commande pour la couleur d'accentuation
\newcommand{\myul}[2][black]{\setulcolor{#1}\ul{#2}\setulcolor{black}}
\newcommand\tab[1][1cm]{\hspace*{#1}}

\usepackage[margin=2cm]{geometry}
\usepackage{eso-pic}         % Pour ajouter des éléments en arrière-plan

\usepackage{enumitem}
%---------------------------------------
% Définir un compteur pour les exemples
\newcounter{exemple}

% Définir la commande \exemple pour afficher un exemple numéroté
\newcommand{\exemple}{%
  \refstepcounter{exemple}% Incrémenter le compteur
  \textbf{\textcolor{orange}{Exemple \theexemple : }} \ignorespaces
}
%---------------------------------------
\newcounter{solution}

% Définir la commande \solutione pour afficher un solution numéroté
\newcommand{\solution}{%
  \refstepcounter{solution}% Incrémenter le compteur
  \textbf{\textcolor{orange}{Solution \thesolution : }} \ignorespaces
}
%---------------------------------------
\definecolor{myorange}{rgb}{1.0, 0.8, 0.0}

% Définir un compteur pour les exercices d'application
\newcounter{exerciceapp}

% Définir la commande pour afficher un exercice d'application numéroté
\newcommand{\exerciceapp}{%
  \refstepcounter{exerciceapp}%
  \textbf{\textcolor{myorange}{Exercice d'application \theexerciceapp :}} \ignorespaces
}
%--------------------------------------
% Définir un compteur pour les exercices d'application
\newcounter{correction}

% Définir la commande pour afficher un correction exercice d'application numéroté
\newcommand{\correction}{%
  \refstepcounter{correction}%
  \textbf{\textcolor{myorange}{Correction \thecorrection :}} \ignorespaces
}
%--------------------------------------
% Définir un compteur pour les remarque d'application
\newcounter{remarque}

%----------------------------------------
\definecolor{myorange1}{rgb}{1.0, 1.5, 0}
% Définir la commande pour afficher un remarque numéroté
\newcommand{\remarque}{%
  \refstepcounter{remarque}%
  \textbf{\textcolor{myorange1}{Remarque \theremarque :}} \ignorespaces
}
% Commande pour ajouter du texte en arrière-plan
\AddToShipoutPicture{
    \AtTextCenter{%
        \makebox[0pt]{\rotatebox{45}{\textcolor[gray]{0.9}{\fontsize{5cm}{5cm}\selectfont Pathé Gobel BA}}}
    }
}

\begin{document}

\maketitle
\newpage
\section*{\underline{\textbf{\textcolor{red}{I. Equations}}}}
\subsection*{\underline{\textbf{\textcolor{red}{1. Compléments sur les trinômes du \( 2^{nd} \)degré}}}}
\begin{itemize}

\item \textcolor{red}{\underline{Forme canonique d'un trinôme du second degré}}

Soit 
\(
P(x)=ax^2+bx+c
\)
où $a,b,c\in\mathbb{R}$ avec $a\neq0$.

Alors
\(
P(x)=a\left[\left(x+\frac{b}{2a}\right)^2-\frac{\Delta}{4a^2}\right]
\)
où
\(
\Delta=b^2-4ac.
\)

(On lit $\Delta$ « delta ».)

\medskip

L’expression obtenue dans $(*)$ est appelée \textcolor{red}{forme canonique} du trinôme du $2^{\text{nd}}$ degré $P(x)$.

\medskip

$\Delta$ est appelé le \textcolor{red}{discriminant} de $P(x)$.

\item \textcolor{red}{\underline{Théorème}}

Soit 
\(
P(x)=ax^2+bx+c
\)
où $a,b,c\in\mathbb{R}$ avec $a\neq0$.

\(
\Delta=b^2-4ac
\)

\begin{enumerate}
\item[(i)] Si $\Delta<0$ alors $P$ n’admet pas de racine.

\item[(ii)] Si $\Delta>0$ alors $P$ admet deux racines distinctes $x_1$ et $x_2$ telles que :

\(
x_1=\frac{-b-\sqrt{\Delta}}{2a}
\quad \text{et} \quad
x_2=\frac{-b+\sqrt{\Delta}}{2a}.
\)

\item[(iii)] Si $\Delta=0$ alors $P$ admet une racine double $x_0$ telle que :

\(
x_0=\frac{-b}{2a}.
\)
\end{enumerate}

\item \textcolor{red}{\underline{Factorisation de $P(x)$}}

\medskip

Soit $P(x)=ax^2+bx+c$, $a,b,c\in\mathbb{R}$ avec $a\neq0$.

\begin{enumerate}
\item[(i)] Si $\Delta<0$ alors $P(x)$ n’a pas de factorisation dans $\mathbb{R}$.

\item[(ii)] Si $\Delta>0$ alors $P(x)$ est factorisable telle que :

\(
\boxed{P(x)=a(x-x_1)(x-x_2)}
\)

\item[(iii)] Si $\Delta=0$ alors $P(x)$ est factorisable telle que :

\(
\boxed{P(x)=a(x-x_0)^2}
\)
\end{enumerate}

\end{itemize}

\subsubsection*{\underline{\textbf{\textcolor{red}{a. Équation bicarrée et équation symétrique}}}}
\begin{itemize}
\item \textcolor{red}{Équation bicarrée}

$3x^{4}-4x^{2}+1=0\quad$ (1)

Posons $X=x^{2}$, donc l'équation (1) devient $3X^{2}-4X+1=0\quad$ (2).

Résolvons l'équation (2). On a :

$
\begin{array}{rcl} 
\Delta=16-12=4 & \Rightarrow & X_{1}=\dfrac{4-2}{6}=\dfrac{1}{3}\quad\text{ et }\quad X_{2}=\dfrac{4+2}{6}=1 \\ \\ 
& \Rightarrow & X_{1}=\dfrac{1}{3}=x^{2}\quad \text{ et }\quad X_{2}=1=x^{2} \\ 
\\ & \Rightarrow & \left\lbrace\begin{array}{lll} x_{1}=\dfrac{\sqrt{3}}{3}\quad \text{ ou }\quad x_{1}=-\dfrac{\sqrt{3}}{3}\\ \\ x_{2}=1\quad \text{ ou }\quad x_{2}=-1
\end{array}\right.
\end{array}
$

$$S=\left\lbrace-1;\ -\dfrac{\sqrt{3}}{3};\ \dfrac{\sqrt{3}}{3};\ 1\right\rbrace$$

\textbf{\exerciceapp}

$x^{4}+4x^{2}-5=0\quad$ (3)

Soit $X=x^{2}$, donc l'équation (3) devient $X^{2}+4X-5=0\quad$ (4).

Résolvons l'équation (4). On a :

$
\begin{array}{rcl}
\Delta=16+20=36 & \Rightarrow & X_{1}=\dfrac{-4-6}{2}=-5\quad\text{ et }\quad X_{2}=\dfrac{-4+6}{2}=1 \\ \\ & \Rightarrow & X_{1}=x^{2}=-5\;\text{ (impossible)}\quad \text{ et }\quad X_{2}=1=x^{2} \\ \\ & \Rightarrow &  x_{2}=1\quad \text{ ou }\quad x_{2}=-1
\end{array}
$

$$S=\left\lbrace-1\;;\ 1\right\rbrace$$

\item \textcolor{red}{Équation symétrique}



\end{itemize}
\subsubsection*{\underline{\textbf{\textcolor{red}{b. Équation paramétrique }}}}

\subsection*{\underline{\textbf{\textcolor{red}{2. Equations irrationnelles}}}}
\subsubsection*{\underline{\textbf{\textcolor{red}{a. Équations du type \( \sqrt{f(x)} = \sqrt{g(x)} \)}}}}

\( \sqrt{x^{2}+3x-1} = \sqrt{x+2} \) est une équation irrationnelle du type \( \sqrt{f(x)} = \sqrt{g(x)} \) avec :

\[ f(x) = x^{2}+3x-1 \quad \text{et} \quad g(x) = x+2. \]

Pour résoudre une équation du type \( \sqrt{f(x)} = \sqrt{g(x)} \), on peut utiliser deux méthodes :

\underline{\textbf{Méthode de Résolution}}

\begin{itemize}
    \item \textbf{1ère étape :} Vérification des conditions de validité.
\end{itemize}

On a les équivalences suivantes :

\[
\sqrt{ f(x)} = \sqrt{ g(x)} \Leftrightarrow \begin{cases}
f(x) \geq 0 \\
g(x) \geq 0 \\
f(x) = g(x)
\end{cases}
\]

**Explication :** Pour que les deux racines soient définies et égales, il faut que \( f(x) \) et \( g(x) \) soient toutes deux positives, d'où les conditions \( f(x) \geq 0 \) et \( g(x) \geq 0 \). Ensuite, il faut que les deux expressions soient égales, soit \( f(x) = g(x) \).

En simplifiant ces conditions, on obtient :

\[
\Leftrightarrow \begin{cases}
g(x) \geq 0 \\
f(x) = g(x)
\end{cases}
\]

**Explication :** Ici, on peut réduire les conditions en ne gardant que \( g(x) \geq 0 \), car l'égalité \( f(x) = g(x) \) implique automatiquement que \( f(x) \geq 0 \) lorsque \( g(x) \geq 0 \).

Finalement, cela se résume à :

\[
\Leftrightarrow \begin{cases}
f(x) \geq 0 \\
f(x) = g(x)
\end{cases}
\]

**Explication :** Cette équivalence montre que, pour que l'équation soit valide, il faut vérifier que \( f(x) \geq 0 \), tout en résolvant l'équation \( f(x) = g(x) \).

 \textbf{\exemple}

Résoudre l'équation:

\[
\sqrt{2x + 1} = \sqrt{x + 3} \quad \text{et} \quad \sqrt{x^2 + 3x - 1} = \sqrt{x + 2}
\]

\textbf{\solution}

\[
\sqrt{2x + 1} = \sqrt{x + 3}
\]


\begin{itemize}
    \item \textbf{Étape 1 : Conditions de validité $D_v$}
    
    Pour que les racines carrées existent, il faut que les expressions à l'intérieur des racines soient positives ou nulles :
    \[
    2x + 1 \geq 0 \implies x \geq -\frac{1}{2}
    \]
    \[
    x + 3 \geq 0 \implies x \geq -3
    \]
    Donc, la condition de validité est :
    \[
    D_v = \left[-\frac{1}{2}; +\infty \right[
    \]
    
    \item \textbf{Étape 2 : Résolution de l'équation}
    
    On élève les deux membres de l'équation au carré :
    \[
    (\sqrt{2x + 1})^2 = (\sqrt{x + 3})^2
    \]
    Ce qui donne :
    \[
    2x + 1 = x + 3
    \]
    
    \item \textbf{Étape 3 : Simplification de l'équation}
    
    Isolons $x$ :
    \[
    2x + 1 = x + 3
    \]
    \[
    2x - x = 3 - 1
    \]
    \[
    x = 2
    \]
    
    \item \textbf{Étape 4 : Vérification de la solution}
    
    Vérifions que $x = 2$ appartient bien à l'ensemble de validité $D_v$ :
    
    \begin{itemize}
        \item $x = 2$ est bien dans $\left[-\frac{1}{2}; +\infty \right[$.
        \item Vérification dans l'équation initiale :
        \[
        \sqrt{2 \times 2 + 1} = \sqrt{2 + 3}
        \]
        \[
        \sqrt{5} = \sqrt{5}
        \]
    \end{itemize}
    
    La solution est donc correcte.
    
\end{itemize}

\textbf{Conclusion :} La solution de l'équation est $x = 2$.

\[
\sqrt{x^2 + 3x - 1} = \sqrt{x + 2}
\]

\begin{itemize}
    \item \textbf{Étape 1 : Conditions de validité $D_v$}
    
    L'équation existe ssi $x^2 + 3x - 1 \geq 0$ et $x + 2 \geq 0$.\\
    La deuxième condition nous donne $x \geq -2$.\\
    Donc $D_v = [-2; +\infty[$.
    
    \item \textbf{Étape 2 : Résolution de l'équation}
    
    Élevons les deux membres de l'équation au carré :
    \[
    x^2 + 3x - 1 = x + 2
    \]
    Cela revient à résoudre l'équation :
    \[
    x^2 + 3x - 1 = x + 2 \implies x^2 + 3x - x - 2 = 0
    \]
    \[
    x^2 + 2x - 3 = 0
    \]
    
    \textbf{Conclusion :} La seule solution est $x = 1$.
\end{itemize}

\begin{itemize}
    \item[•] \textbf{Étape 1 : Conditions de validité}
\end{itemize}

On doit vérifier que \( f(x) = x^2 + 3x - 1 \geq 0 \) et \( g(x) = x + 2 \geq 0 \). La deuxième condition nous donne \( x \geq -2 \).

\begin{itemize}
    \item[•] \textbf{Étape 2 : Résolution de l'équation}
\end{itemize}

Élevons les deux membres de l'équation au carré :

\[
x^2 + 3x - 1 = x + 2
\]

Cela revient à résoudre l'équation :

\[
x^2 + 3x - 1 = x + 2 \quad \Leftrightarrow \quad x^2 + 2x - 3 = 0
\]

\begin{itemize}
    \item[•] \textbf{Étape 3 : Résolution de l'équation quadratique}
\end{itemize}

On résout \( x^2 + 2x - 3 = 0 \) par factorisation :

\[
(x - 1)(x + 3) = 0 \quad \Rightarrow \quad x = 1 \quad \text{ou} \quad x = -3
\]

\begin{itemize}
    \item[•] \textbf{Étape 4 : Vérification des solutions}
\end{itemize}

- Pour \( x = 1 \) :

\[
f(1) = 1^2 + 3(1) - 1 = 3 \quad \text{et} \quad g(1) = 1 + 2 = 3
\]

Les deux fonctions sont positives et égales, donc \( x = 1 \) est une solution.

- Pour \( x = -3 \) :

\[
f(-3) = (-3)^2 + 3(-3) - 1 = 9 - 9 - 1 = -1 \quad \text{et} \quad g(-3) = -3 + 2 = -1
\]

Les deux fonctions sont négatives, donc \( x = -3 \) n'est pas une solution.

\textbf{Conclusion :} La seule solution est \( x = 1 \).

\subsubsection*{\underline{\textbf{\textcolor{red}{b. Équations du type \( \sqrt{f(x)} = ax + b \)}}}}

\underline{\textbf{Méthode de Résolution}}

\[\text{On a les équivalences suivantes :}\]

\[
\sqrt{f(x)} = g(x) \Leftrightarrow \begin{cases} 
f(x) \geq 0 \\
g(x) \geq 0 \\
f(x) = g(x)^{2} 
\end{cases}
\]

**Explication :**
  
- \( f(x) \geq 0 \) : Pour que la racine carrée \( \sqrt{f(x)} \) soit définie, il est nécessaire que l'expression sous la racine soit positive ou nulle, d'où \( f(x) \geq 0 \).

- \( g(x) \geq 0 \) : Comme la racine carrée d'un nombre est toujours positive, il faut également que \( g(x) \), qui représente l'autre membre de l'équation, soit positif ou nul.

- \( f(x) = g(x)^2 \) : Une fois ces conditions vérifiées, on élève les deux membres de l'équation au carré afin d'éliminer la racine carrée, ce qui donne \( f(x) = g(x)^2 \).

En simplifiant, on obtient :

\[
\Leftrightarrow \begin{cases}
g(x) \geq 0 \\
f(x) = g(x)^{2}
\end{cases}
\]

**Explication :**
  
Dans certains cas, la condition \( f(x) \geq 0 \) devient redondante, car l'égalité \( f(x) = g(x)^2 \) implique que \( f(x) \) est nécessairement positive si \( g(x) \geq 0 \).

 \textbf{\exemple}

Résoudre l'équation:

\( \sqrt{2x^{2}+4x-6} = 2x - 1 \)

\textbf{\solution}

\textbf{Étape 1 : Conditions de validité $D_v$}

\[
2x-1\geq 0
\]

\[
\Longleftrightarrow x\geq \frac{1}{2}.
\]


Sous cette condition, on peut élever les deux membres au carré :

\textbf{Étape 2 : Résolution de l'équation }
\[
\begin{aligned}
2x^{2}+4x-6&=(2x-1)^2\\
2x^{2}+4x-6&=4x^{2}-4x+1\\
2x^{2}-8x+7&=0.
\end{aligned}
\]

Calculons le discriminant :

\[
\Delta=(-8)^2-4(2)(7)=64-56=8.
\]

Ainsi,

\[
\begin{aligned}
x_1&=\frac{8-\sqrt{8}}{4}
=\frac{8-2\sqrt{2}}{4}
=2-\frac{\sqrt{2}}{2},\\[2mm]
x_2&=\frac{8+\sqrt{8}}{4}
=\frac{8+2\sqrt{2}}{4}
=2+\frac{\sqrt{2}}{2}.
\end{aligned}
\]

\textbf{Étape 3 : Vérification de l'appartenance des solution dans $D_v$ }

Les deux solutions vérifient la condition \(x\geq\frac12\).

\[
\boxed{
S=
\left\{
2-\frac{\sqrt{2}}{2}\,;\,
2+\frac{\sqrt{2}}{2}
\right\}
}
\]

\section*{\underline{\textbf{\textcolor{red}{II. Inéquations irrationnelles}}}}
\subsection*{\underline{\textbf{\textcolor{red}{1. Inéquations du type \( \sqrt{f(x)} \leq ax+b \) }}}}

\underline{\textbf{Méthode de Résolution}}

\[\text{On a les équivalences suivantes :}\]

\[
\sqrt{f(x)} \leq g(x) \Leftrightarrow \begin{cases} 
f(x) \geq 0 \\
g(x) \geq 0 \\
f(x) \leq g(x)^{2}
\end{cases}
\]

**Explication :**  

- **\( f(x) \geq 0 \)** : L'expression sous la racine carrée doit être positive ou nulle pour que la racine soit définie.

- **\( g(x) \geq 0 \)** : Comme la racine carrée est toujours positive ou nulle, il est nécessaire que l'expression \( g(x) \), ici \( ax + b \), soit aussi positive ou nulle.

- **\( f(x) \leq g(x)^2 \)** : On élève les deux membres de l'inéquation au carré pour éliminer la racine, ce qui donne l'inéquation \( f(x) \leq g(x)^2 \).


\textbf{\exemple}

Résolvons l'inéquation \( \sqrt{2x+1} \leq x+3 \).

\textbf{\solution}

\begin{itemize}
    \item \textbf{Étape 1 : Conditions de validité \( D_v \)}

	$
	\begin{aligned}
		\begin{cases}
	2x + 1 \geq 0\\
	x + 3 \geq 0
	\end{cases} &\implies
	\begin{cases}
	x \geq -\dfrac{1}{2}\\
	x \geq -3
	\end{cases}\\ &\implies
		\begin{cases}
	x \geq -\dfrac{1}{2}\\
	x \geq -3
	\end{cases}\\ &\implies
		\begin{cases}
	\left[-\dfrac{1}{2}; +\infty \right[\\[3mm]
	\left[-3; +\infty \right[
	\end{cases}
	\end{aligned}
	$    

$D_v=\left[-\dfrac{1}{2}; +\infty \right[ \cap \left[-3; +\infty \right[$

    \item \textbf{Étape 2 : Résolution de l'inéquation}
    
    On élève les deux membres de l'inéquation au carré :

                                    $
                                    \begin{aligned}
                                    2x + 1 \leq (x + 3)^2 \quad &\Leftrightarrow \quad 2x + 1 \leq x^2 + 6x + 9\\
                                    \quad &\Leftrightarrow \quad x^2 + 4x + 8 \geq 0
                                    \end{aligned}
                                    $

Posons $x^2+4x+8=0$

$\Delta=16-32=-16$

Comme $\Delta<0,$ donc le signe de $x^2+4x+8$ du signe 1 qui est positif donc \( \mathbb{R} \)


\textbf{Conclusion :} L'ensemble des solutions est 

                                        $
                                        \begin{aligned}
                                        S&=D_v \cap \mathbb{R}\\
                                        &=\left[-3; +\infty \right[ \cap \mathbb{R}\\
                                        &=\left[-3; +\infty \right[
                                        \end{aligned}
                                        $
                                        
$\boxed{ S=\left[-3; +\infty \right[}$

\end{itemize}


\textbf{\exemple}

Résolvons l'inéquation \( \sqrt{x^{2}+3x-1} \leq x+2 \).

\textbf{\solution}

\begin{itemize}
    \item \textbf{Étape 1 : Conditions de validité \( D_v \)}
\end{itemize}

L'inéquation est définie si \( x^{2} + 3x - 1 \geq 0 \) et \( x + 2 \geq 0 \).

- \( x + 2 \geq 0 \) donne \( x \geq -2 \).

- Pour \( x^{2} + 3x - 1 \geq 0 \), on trouve les racines du polynôme : \( x = -3.41 \) et \( x = 0.41 \).

Le domaine de validité est donc \( D_v = [-2; +\infty[ \).

\begin{itemize}
    \item \textbf{Étape 2 : Résolution de l'inéquation}
\end{itemize}

On élève les deux membres de l'inéquation au carré :

\[
x^{2} + 3x - 1 \leq (x + 2)^{2}
\]

Cela revient à résoudre l'inéquation :

\[
x^{2} + 3x - 1 \leq (x + 2)^{2} \quad \Leftrightarrow \quad x^{2} + 3x - 1 \leq x^{2} + 4x + 4.
\]

En simplifiant, on obtient :

\[
0 \leq x + 5.
\]

Cette inéquation est vérifiée pour tout \( x \geq -5 \).

\textbf{Conclusion :} L'ensemble des solutions est \( S = [-2; +\infty[ \), car il faut tenir compte des conditions initiales.

\subsection*{\underline{\textbf{\textcolor{red}{2. Inéquations du type \( \sqrt{f(x)} \geq ax+b \) }}}}

\underline{\textbf{Méthode de Résolution}}

\[
\text{On a les équivalences suivantes :}
\]

\[
\begin{aligned}
\sqrt{f(x)} \geq g(x)
&\Leftrightarrow
\begin{cases}
f(x) \geq 0\\
g(x) \geq 0\\
f(x) \geq g(x)^2
\end{cases}
\quad \textbf{ou} \quad
\begin{cases}
f(x) \geq 0\\
g(x) \leq 0
\end{cases}\\
&\Leftrightarrow
\boxed{
\begin{cases}
g(x) \geq 0\\
f(x) \geq g(x)^2
\end{cases}
\quad \textbf{ou} \quad
\begin{cases}
f(x) \geq 0\\
g(x) \leq 0
\end{cases}}
\end{aligned}
\]

\textbf{Explication :}

\begin{itemize}
    \item \textbf{\( f(x) \geq 0 \)} : Comme la racine carrée
    \( \sqrt{f(x)} \) doit être définie, il est nécessaire que
    l'expression sous la racine soit positive ou nulle.

    \item \textbf{Premier cas : \( g(x) \geq 0 \)} :
    Les deux membres sont alors positifs ou nuls. On peut donc
    élever les deux membres au carré sans changer le sens de
    l'inégalité :
    \[
    \sqrt{f(x)} \geq g(x)
    \Leftrightarrow
    f(x) \geq g(x)^2.
    \]

    \item \textbf{Deuxième cas : \( g(x) \leq 0 \)} :
    Comme une racine carrée est toujours positive ou nulle, on a
    automatiquement
    \[
    \sqrt{f(x)} \geq 0 \geq g(x).
    \]
    Ainsi, dans ce cas, il suffit d'avoir
    \[
    f(x) \geq 0.
    \]
\end{itemize}

\textbf{En simplifiant, on obtient :}

\[
\boxed{
\sqrt{f(x)} \geq g(x)
\Leftrightarrow
\left[
\begin{array}{l}
g(x)\geq0\\
f(x)\geq g(x)^2
\end{array}
\right]
\quad\textbf{ou}\quad
\left[
\begin{array}{l}
f(x)\geq0\\
g(x)\leq0
\end{array}
\right]
}
\]

\textbf{Dans le cas particulier où \(g(x)=ax+b\), on résout donc
séparément les deux systèmes :}

\[
\begin{cases}
f(x)\geq0\\
ax+b\leq0
\end{cases}
\]

et

\[
\begin{cases}
ax+b\geq0\\
f(x)\geq(ax+b)^2.
\end{cases}
\]

\textbf{L'ensemble solution de l'inéquation est l'union des
solutions de ces deux systèmes.}


\textbf{\exemple}

Résolvons l'inéquation \( \sqrt{2x+1} \geq x+3 \).

\textbf{\solution}


   \[
                                    \textbf{Posons }
                                    S_{1}:
                                    \begin{cases}
                                    2x+1 \geq 0\\[3mm]
                                    x+3 \leq 0
                                    \end{cases}\qquad
                                    \textbf{Posons }
                                    S_{2}:
                                    \begin{cases}
                                    x+3 \geq 0\\[3mm]
                                    2x+1 \geq (x+3)^2
                                    \end{cases}
                                    \]

Ce qui correspond à :

\[
S_{1}:
\begin{cases}
2x + 1 \geq 0 \\ \\
x + 3 \leq 0 
\end{cases}
\textbf{ et }
S_{2}:
\begin{cases} 
x + 3 \geq 0 \\[3mm]
x^2 + 4x + 8 \leq 0
\end{cases}
\]

\[
\textbf{Posons }
S_{1}:
\begin{cases}
2x + 1 = 0 \\ \\
x + 3 = 0 
\end{cases}
\textbf{Posons }
S_{2}:
\begin{cases} 
x + 3 = 0 \\[3mm]
x^2 + 4x + 8 = 0
\end{cases}
\]

\underline{Résolvons \( S_1 \)} :

\(
\begin{cases}
x = \dfrac{-1}{2} \\
x =  -3 
\end{cases} 
\)


    \begin{tikzpicture}[node style/.style={fill opacity=0,text opacity=1}]
        \tkzTabInit[espcl=1.75]{$x$/.5,$2x + 1$/.7,$ x + 3 $/.7}{$-\infty$,$-3$,$\frac{-1}{2}$,$+\infty$}
        \tkzTabLine{,-,t,-,z,+, }
        \tkzTabLine{,-,z,+,t,+, }
    \end{tikzpicture}
    
$S_1 = \varnothing$


\underline{Résolvons \( S_2 \)} :

\(
\begin{cases} 
x + 3 = 0 \\[3mm]
x^2 + 4x + 8 = 0
\end{cases}
\implies
\begin{cases} 
x = -3 \\[3mm]
\text{Pas de Racine }
\end{cases}
\)

    \begin{tikzpicture}[node style/.style={fill opacity=0,text opacity=1}]
        \tkzTabInit[espcl=2,lgt=3]{$x$/.5,$x + 3$/.7,$ x^2 + 4x + 8 $/.7}{$-\infty$,$-3$ ,$+\infty$}
        \tkzTabLine{ ,-,z,+, }
        \tkzTabLine{ ,+,t,+, }
    \end{tikzpicture}

$S_2 = \varnothing$

\subsection*{\underline{\textbf{\textcolor{red}{3. Inéquations du type \( \sqrt{f(x)} \leq \sqrt{g(x)} \)}}}}


\underline{\textbf{Méthode de résolution}}

On a : $\sqrt{f(x)}\leq\sqrt{g(x)}\;\Leftrightarrow\;\left\lbrace\begin{array}{lll} f(x) \geq 0 \\ g(x)\geq 0\\ f(x)\leq g(x)\end{array}\right.$

\textbf{\exemple}

Résoudre dans \(\mathbb{R}\) l'inéquation :


\[
\sqrt{2x+3}\leq\sqrt{x+5}.
\]

\textbf{\solution}

\textbf{1. Conditions d'existence}

\[
\left\{
\begin{array}{l}
2x+3\geq0\\
x+5\geq0
\end{array}
\right.
\]

soit

\[
\left\{
\begin{array}{l}
x\geq-\dfrac{3}{2}\\
x\geq-5.
\end{array}
\right.
\]

Ainsi,

\[
D=\left[-\frac{3}{2},+\infty\right[.
\]

\textbf{2. Résolution de l'inéquation}

Sur \(D\), on peut comparer directement les radicandes :

\[
\sqrt{2x+3}\leq\sqrt{x+5}
\Longleftrightarrow
2x+3\leq x+5.
\]

Donc :

\[
x\leq2.
\]

\textbf{3. Intersection}

On obtient :

\[
x\in
\left[-\frac{3}{2},+\infty\right[
\cap
]-\infty,2].
\]

Par conséquent :

\[
\boxed{
S=\left[-\frac{3}{2},2\right]
}
\]

\textbf{\exemple}

Résoudre dans $\mathbb{R}\;\sqrt{x^{2}-4x}\leq\sqrt{x+7}$.

\textbf{\solution}

$\begin{array}{rcl} \sqrt{x^{2}-4x}\leq\sqrt{x+7} & \Leftrightarrow & \left\lbrace\begin{array}{rcl} x^{2}-4x&\geq&0 \\ x+7&\geq&0 \\ (\sqrt{x^{2}-4x})^{2}&\leq&(\sqrt{x+7})^{2} \end{array}\right. \\ \\ & \Leftrightarrow & \left\lbrace\begin{array}{rcl} x(x-4)&\geq&0 \\ x&\geq&-7 \\ x^{2}-4x&\leq&x+7 \end{array}\right. \\ \\ & \Leftrightarrow & \left\lbrace\begin{array}{lll} x\in\;(]-\infty\;;\ 0]\cup[4\;;\ +\infty[)\cap[-7\;;\ +\infty[ \\ x^{2}-5x-7\leq 0 \end{array}\right. \\ \\  & \Leftrightarrow & \left\lbrace\begin{array}{lll} x\in\;D_{E}=]-7\;;\ 0]\cup[4\;;\ +\infty[ \\ x^{2}-5x-7\leq 0 \end{array}\right. \end{array}$

Soit l'équation $x^{2}-5x-7=0$

On a : $\Delta=25+28=53\;\Rightarrow\;\sqrt{\Delta}=\sqrt{53}.$

Soient $x_{1}$ et $x_{2}$ les solutions distinctes de $x^{2}-5x-7=0$.

Alors, $x_{1}=\dfrac{5-\sqrt{53}}{2}\ $ et $\ x_{2}=\dfrac{5+\sqrt{53}}{2}$

Ainsi, sans aucune contrainte, l'inéquation $x^{2}-5x-7\leq 0$ admettra comme solution $S_{1}=[x_{1}\;;\ x_{2}]$

$$\begin{array}{|c|lclclclclcr|} \hline x & -7 & & x_{1} &  & 0 &  & 4 & & x_{2} & & +\infty \\ \hline x^{2}-4x & & + & | & + & 0 & - & 0 & + &|& + & \\ \hline x^{2}-5x-7 &  & + & 0 & - & | & - & | & - & 0 & + & \\ \hline S &  & + & 0 & (-) & | & + & |& (-) & 0 & + & \\ \hline \end{array}$$

Donc en tenant compte des conditions initiales, l'ensemble des solutions du problème sera donné par $$S=S_{1}\cap D_{E}=[x_{1}\;;\ 0]\cup[4\;;\ x_{2}]$$

$\centerdot\ \ \sqrt{f(x)}\geq g(x)$

On a : $\sqrt{f(x)}\geq g(x)\;\Leftrightarrow\;\left\lbrace\begin{array}{lll} g(x)\geq 0\\ f(x)\geq (g(x))^{2} \end{array}\right.\quad\text{ou}\quad\left\lbrace\begin{array}{lll} f(x)\geq 0 \\ g(x)&lt; 0\end{array}\right.$

Donc, on obtient deux solutions $S_{1}$ et $S_{2}.$

Ainsi, la solution finale est donnée par $S=S_{1}\cup S_{2}.$


\textbf{\exerciceapp}

Résoudre dans \( \mathbb{R} \)

\begin{itemize}
\item \( \sqrt{x+3} \leq 5 \)
\item \( \sqrt{-2x+1} \geq -3 \)
\item \( \sqrt{x+3} \leq \sqrt{-2x+4} \)
\item \( \sqrt{x^{2}-4} \leq x+1 \)
\item \( \sqrt{x^{2}+2x -3} \leq x-2 \)
\item \( \sqrt{x^{2}-6x}=\sqrt{x-6} \)
\end{itemize}
\textbf{\correction}
\section*{\underline{\textbf{\textcolor{red}{III. Systèmes}}}}
\subsection*{\underline{\textbf{\textcolor{red}{1. Systèmes de 3 équations linéaires à trois inconnues}}}}

Le système 
\[
\begin{cases}
a_1x + b_1y + c_1z = d_1 \\
a_2x + b_2y + c_2z = d_2 \\
a_3x + b_3y + c_3z = d_3
\end{cases}
\]
est formé de trois équations. Chacune d’elle contient trois inconnues \(x\), \(y\) et \(z\) avec des exposants tous égaux à 1. On dit qu’on a un système de trois équations linéaires à trois inconnues \(x\), \(y\) et \(z\).

Résoudre un tel système c’est trouver tous les triplets \((x, y, z)\) de nombres réels qui vérifient les trois équations du système.

\subsubsection*{\underline{\textbf{\textcolor{red}{ Exemple }}}}

Exemple de système :

\[
\begin{cases}
2x + y - z = 3 \\
x - 2y + 3z = -4 \\
3x + y + 2z = 7
\end{cases}
\]

\subsubsection*{\underline{\textbf{\textcolor{red}{Exemple de système triangulaire}}}}

Un système triangulaire supérieur est un système d'équations où chaque équation contient une inconnue de moins que la précédente, formant un triangle.

$
\left\{
\begin{array}{rcc l}
2x + 3y - z & = & 5 & (L_1) \\[2mm]
y + 4z & = & -2 & (L_2) \\[2mm]
y + z & = & 3 & (L_3)
\end{array}
\right.
$

Dans ce système, la dernière équation contient seulement \( z \), ce qui permet de le résoudre directement, puis de remonter en substituant dans les équations précédentes pour obtenir \( y \) et \( x \).

\subsubsection*{\underline{\textbf{\textcolor{red}{b. Méthode du pivot de Gauss }}}}

a. Exemple 1

	Résolvons le système suivant en utilisant la méthode du pivot de Gauss

	$\left\lbrace\begin{array}{lcl}x+10y-3z& =&5\\2x-y+2z& = 2\\-x+y+z& =& -3\end{array}\right.$

	Pour résoudre un tel système, on commence par désigner les $3$ équations respectivement par $L_{1}$ ; $L_{2}$ et $L_{3}$

	$\left\lbrace\begin{array}{lcl}x+10y-3z& =& 5\qquad\left(L_{1}\right)\\2x-y+2z& =& 2\qquad\left(L_{2}\right)\\-x+y+z& =& -3\qquad\left(L_{3}\right)\end{array}\right.$

	$1^{er}$ étape : 


	On fixe l'équation $\left(L_{1}\right)$ puis on élimine l'inconnue $x$ dans $\left(L_{2}\right)$ en considérant le sous-système $-2\left\lbrace\begin{array}{lcl}x+10y-3z& =& 5\qquad\left(L_{1}\right)\\2x-y+2z& =& 2\qquad\left(L_{2}\right)\end{array}\right.$


	$$\text{d'où }\left\lbrace\begin{array}{lcl}-2x-20y+6z& =& -10\\2x-y+2z& =& 2\end{array}\right.$$
	
	ainsi on a : $-21y+8z=-8\quad\left(L'_{2}\right)$

$2^{ième}$ étape :

	On fixe l'équation $\left(L_{1}\right)$ puis on élimine l'inconnue $x$ dans $\left(L_{3}\right)$ en considérant le sous-système $\left\lbrace\begin{array}{lcl}x+10y-3z&=&5\qquad\left(L_{1}\right)\\-x+y+z&=&-3\qquad\left(L_{3}\right)\end{array}\right.$

	d'où ainsi on a : $11y-2z=2\quad\left(L'_{3}\right)$

	Ainsi, on obtient le système équivalent suivant : 
	$
    \left\{
    \begin{array}{rcc l}
    x + 10y - 3z &=& 5 & (L_1) \\
    -21y + 8z &=& -8 & (L'_2) \\
    11y - 2z &=& 2 & (L'_3)
    \end{array}
    \right.
    $

$3^{ième}$ étape :

	On fixe $\left(L'_{2}\right)$ puis on élimine $y$ dans $\left(L'_{3}\right)$ en considérant le sous-système $\left\lbrace\begin{array}{lcl}-21y+8z& =& -8\qquad\left(L'_{2}\right)\\11y-2z& =& 2\qquad\left(L'_{3}\right)\end{array}\right.$

	Ainsi on a $46z=-46\left(L"_{3}\right)$

	On obtient le système équivalent suivant dit système triangulaire

												$ \left\{
                                                \begin{array}{rcc l}
                                                x + 10y - 3z & = & 5 & (L_1) \\
                                                -21y + 8z & = & -8 & (L'_2) \\
                                                46z & = & -46 & (L''_3)
                                                \end{array}
                                                \right.$
	
$4^{ième}$ étape :

	Pour terminer la résolution, on détermine $z$ dans $\left(L"_{3}\right)$, puis on remplace $z$ par cette valeur dans $\left(L'_{2}\right)$, on obtient la valeur de $y$ puis on obtient celle de $x$, en substituant à $y$ et $z$ par leurs valeurs respectives dans $\left(L'_{1}\right).$

	Ainsi on a : $S=\left\lbrace(2\ ;\ 0\ ;\ -1)\right\rbrace$
\end{document}
